\subsubsection{Tiền xử lý}

\paragraph{Xử lý các giá trị còn thiếu} Bộ dữ liệu California Housing không có
giá trị thiếu nên bước này được bỏ qua.

\paragraph{Chuẩn hóa dữ liệu} Các biến đặc trưng được chuẩn hóa về trung bình
$\mu = 0$ và phương sai $\sigma^2 = 1$. Điều này giúp các giá trị của các biến
đặc trưng được đưa về cùng thang đo, giúp quá trình học bằng Gradient Descent
của các mô hình nhanh hơn (các đường mức không còn bị kéo dãn). Ngoài ra, việc
chuẩn hóa như vậy giúp cho việc tính toán không bị tràn số (đã thử và bị với mô
hình Mini-batch Gradient Descent). Lợi ích của cách chuẩn hóa này so với chuẩn
hóa về $[0, 1]$ đó là sự xuất hiện các ngoại lai có giá trị rất lớn làm giá trị
sau khi chuẩn hóa rất nhỏ. Điều này không làm các biến về cùng đúng thang đo so
với mong muốn.

\paragraph{Chia dữ liệu} Dữ liệu được chia thành 3 tập là huấn luyện (train),
xác thực (validation) và kiểm tra (test) với tỉ lệ tương ứng là $0.6$,~$0.2$
và~$0.2$.
